\chapter{HMC Tuning Interfaces}
\label{ch:bf-hmc-tuning-interfaces}

HMC tuning determines which frozen kernels are suitable for further sampling.
For each declared trajectory length $L$, BayesFilter measures proposed step sizes
$\epsilon$, independently verifies the surviving pairs, and retains every pair
that passes. A failed pair can motivate another measured proposal. A passing
pair does not end the search or establish posterior convergence.

\section{One candidate-set procedure}

The public functions \path{tune_hmc_kernel} and
\path{tune_fixed_transport_hmc_kernel} share
\path{HMCTuningCandidateSetController}. Both return
\path{HMCTypedCandidateSetRun}, whose \path{result} contains the full
\path{HMCTuningCandidateSetResult}. Their preparation differs with the target
and coordinates; their candidate scheduling and retention do not.

Ordinary \path{HMCKernelTuningConfig} prepares a windowed mass matrix or an
explicit identity metric. \path{FixedTransportHMCKernelTuningConfig} prepares
identity mass in the coordinates of a reconstructable frozen transport.
A caller with already prepared geometry supplies \path{HMCControllerConfig}
and the numerical binding's \path{typed_adapter}. The conditional
\path{TensorFlowHMCKernelTuningConfig} proposal-field route also uses the shared
scheduler, but retains its weaker mechanics authority.

The capability registry \path{HMC_TUNING_INTERFACE_CAPABILITIES} identifies the
supported target and coordinate combinations. Chain runners, stage helpers,
discovery routines, and historical selectors cannot independently issue new
candidate-set tuning evidence. Unsupported legacy runner customizations fail
before execution. Historical single-kernel results remain readable, with their
original identities; they are not the result of the current procedure.

\lstinputlisting[language=Python,caption={Inspecting supported target and coordinate routes.}]{examples/hmc_tuning_route_selection.py}

\section{Preparation and immutable geometry}

Geometry is a shared responsibility rather than an unspecified tuner default.
The caller supplies the initial center and may supply a local geometry
hypothesis; the tuner validates that hypothesis and owns the subsequent affine
construction and adaptation.  Under windowed adaptation, it tries a supplied
negative log-posterior Hessian, an initial covariance, and positive diagonal
parameter scales, in that order, before falling back to identity.  Thus a
negative Hessian means

\begin{equation}
  H(\theta_0)=-\nabla_{\theta}^{2}\ell(\theta_0),
  \qquad
  \Sigma_0 \approx H(\theta_0)^{-1},
  \label{eq:bf-ordinary-hmc-initial-geometry}
\end{equation}

\noindent
in the same unconstrained coordinates as the supplied center.  With geometry
fallback enabled, a rejected higher-priority hypothesis is recorded before the
next is tried; with fallback disabled, it is a hard error.  The explicit fixed-
identity policy ignores all hints.  Identity is a useful mechanics baseline,
not evidence that an anisotropic posterior has identity geometry.

When no hint is supplied, the initializer uses unit curvature frequencies;
it does not estimate the target Hessian. A finite density or a near-zero score
at the initial point therefore does not establish a stable initial step.
Strongly scaled targets can fail the bootstrap before windowed adaptation or
candidate search begins. Inspect that preparation failure and supply checked
geometry in the active coordinates when available. Increasing the subsequent
candidate budget does not repair a bootstrap that never reaches its handoff.

Bootstrap epsilon decisions use the mean Metropolis probability over all
recorded proposals, including rejected proposals and excluding discarded burnin:
\[
  \bar\alpha=\frac{1}{n}\sum_{t=1}^{n}
     \exp\{\min(\log r_t,0)\}.
\]
The realized accepted fraction includes the additional uniform accept/reject
draw and remains a separate movement diagnostic. For example, a recorded mean
probability of $0.6846$ lies within the configured $[0.65,0.75]$ band even when
$13/16=0.8125$ proposals were accepted. That binary fraction must not cause an
upward epsilon repair. Missing, nonfinite or incomplete probability traces stop
the screen. The ordinary fixed-kernel band and the optional startup floor have
different preparation roles; neither establishes final candidate or posterior
validity.

A target-domain exception permits a bounded smaller bootstrap trial only if
the adapter specifically classifies it and the repository failure record
locates it in a proposal target callback after a finite pre-transition state.
Initial-state, retained-state, unattributed and unrelated execution failures
remain terminal. The failed round retains its original assertion and first
failure; it receives no invented acceptance probability. A smaller finite
parent and the failed epsilon can nominate a logarithmic midpoint; otherwise
the configured multiplier reduces epsilon. The next trial recomputes $L$,
records any clamp and consumes the existing repair budget. A failed reusable
runner is replaced, and only a subsequent completed screen can supply the
repaired preparation pair. Target assertions and Metropolis correction remain
part of every trial.

An optional finite startup probe is available through
\path{HMCKernelTuningConfig(bootstrap_initialization_rounds=20)}. Zero is the
default; a positive integer bounds the number of rounds. Each round tests four
independent momentum proposals and shrinks the step when proposals are invalid
or acceptance is below the configured lower repair floor. The probe preserves
the target, affine mass, failed proposals and separate seed stream. Its selected
pair must still pass fresh bootstrap and operational adaptation. In this mode,
bootstrap uses its \path{warmup_startup_only} role: high acceptance does not
trigger upward refinement before adaptation, and the starting step is not
declared a stability ceiling.

Probe wall times remain in the geometry's full artifact hash for audit, but are
excluded from the separate numerical geometry hash used by new candidate scopes.
Otherwise, two identical preparations could receive different candidate seeds
merely because a probe took longer. Numerical probe outcomes, seed provenance,
mass and starts remain bound to the numerical identity. Existing checkpoints
retain their saved scopes and streams.

The probe follows the declared target-status policy. With \path{none}, it checks
finite retained and proposed states, values and scores without manufacturing
status data. With \path{per_chain_step}, complete valid telemetry is also
required. Invalid initial or retained states stop preparation. A finite startup
nomination neither qualifies a final candidate nor establishes convergence.

Preparation success need not mean that the mass changed. The
\path{allow_valid_incumbent} policy permits zero empirical updates; inspect
\path{operational_metric_update_count} when the initial geometry is poorly
scaled. A regression diagnosis retained identity geometry throughout preparation
and subsequently failed posterior equilibration despite verified tuning
candidates. Ordinary operational preparation runs one chain; it constructs the
four-chain bank for candidate verification afterwards. Its metric assessment
preserves temporal order. For explicitly supplied multiple chains, temporal
information is estimated separately within each chain. Split R-hat is reported
when available and never vetoes or delays a metric proposal.

The proposed position covariance is the finite-window centered second moment,
divided by $N-1$, with correlations shrunk toward zero while preserving marginal
variances. It is a geometry hypothesis during adaptation, not an unbiased
posterior covariance estimate. State-count and temporal-information floors,
numerical rank, condition and regularization-discrepancy checks screen this
hypothesis. Positive-definiteness, coordinate reconstruction, target value/score
parity and a fresh reasonable-step check are required before using its transform.

Ordinary preparation records each completed window's metric decision and failed
checks in \path{preparation_progress.json}. It also reports whether the scheduled
windows can meet the state-count floors. The \path{standard} preset is a local
diagnostic allocation: with the default metric policy, its 150 transitions in six dimensions contain slow
windows of 30, 60 and 30, each shorter than the dense minimum of 64. The explicit
\path{serious} preset funds a larger preparation. Inspect its actual decisions
as well; neither a larger allocation nor an accepted metric update establishes
posterior equilibration. Rejections preserve the incumbent and identify whether
evidence, transform construction, affine parity or step qualification failed.

At a metric change, the fresh step search starts from the previous coordinate
system's bounded dual-averaging average. Clipping in log space before
exponentiation prevents a large unconstrained adaptation proposal from
exhausting the search before any usable step is tested. This old bound limits
only the initial hypothesis: the search may expand in the new coordinates and
must independently qualify its step. Rejected probes preserve their starting
value and attempts. A nonfinite warmup trace identifies its failed fields and
first indices; the failure still vetoes preparation.

The optional \path{metric_evidence_policy="finite_window"} keeps the same
covariance formula and numerical checks while treating temporal ESS as
explanatory. The default \path{temporal_information} retains the ESS floor.
With the optional policy, automatic preparation allocates slow windows large
enough to reach the existing dense state-count floor when the declared budget
permits, or the diagonal floor when only that is affordable. It preserves the
total transitions and initial/final buffers, merging an undersized final slow
remainder. Insufficient total capacity remains visible. This schedule change
removes avoidable count failures; it does not establish adequate covariance
information or posterior equilibration. Directly supplied schedules and the
default policy retain their declared window lengths.
This option permits an early geometry proposal before the chain has explored
well enough to estimate posterior covariance precisely. For any fixed
nonsingular affine factor $A$, $\theta=c+Az$ has transformed score
$A^{\mathsf T}\nabla_\theta\ell$; exact-score Metropolis-adjusted sampling in
the transformed coordinates does not require $AA^{\mathsf T}$ to equal the
posterior covariance. Its usefulness still depends on exploration and numerical
health. Adaptation must finish before candidate verification and retained draws.

The option remains experimental. Its first matched regression preparation
applied two updates and later failed a nonfinite warmup trace. The existing
reasonable-step probe checks a few momenta at one state and cannot guarantee
stability everywhere visited later. Preparation failures remain failures,
and neither policy turns metric qualification into posterior convergence.

The optional \path{metric_probe_num_results} extends each of the four independent
metric-boundary probes to a short chain at fixed epsilon and fixed leapfrog
count. Its default is one transition. A larger allocation, for example sixteen,
tests evolving positions and checks every proposed and retained endpoint,
score, momentum and energy correction before the existing acceptance bracket
can qualify the step. The evidence records the probe length. Retained-state
inconsistency or unclassified execution failure stops preparation; a failed
proposal rejects that probe. This remains an experimental preparation option:
short chains cannot guarantee global stability. Later warmup windows, candidate
measurement and independent verification still require their own numerical
checks. It neither adds R-hat admission nor retries an invalid warmup window.

The separate optional \path{preparation_max_restarts} bounds recovery from
rejected nonfinite proposals; its default is zero. Recovery requires finite
retained states, values and scores, consistent acceptance decisions and valid
declared telemetry. The entire failed preparation attempt is archived and its
adaptation statistics discarded. A full new schedule starts from the failed
window's independently validated starting state and qualified transform, with
fresh streams and reset covariance and dual-averaging statistics. A new probe
may only shrink from half the consumed failing step. High acceptance can
nominate this conservative preparation step; it does not relax final candidate
acceptance or independent verification. A changed metric still requires fresh
qualification in its new coordinates. All discarded work and probe work counts
toward finite attempt and execution budgets. Retained or shared invalidity and
unclassified execution errors remain fatal. This ordinary-preparation option
is experimental: three restarts is a development allocation, not a universal
robustness guarantee.

\lstinputlisting[language=Python,caption={Binding a local covariance estimate in the target coordinates.}]{examples/hmc_tuning_covariance_first.py}


\path{prepare_operational_windowed_mass_handoff} freezes the operational metric
and actual post-warmup start bank. The numerical factory
\path{bind_hmc_candidate_set_execution_from_preparation} preserves both its
bootstrap and final affine layers. The explicit
\path{bind_hmc_candidate_set_execution} factory instead accepts a frozen mass
or supported frozen transport payload. A bound scope records target, data/prior
lineage, source dependencies, coordinates, metric, starts, warmup, dtype,
backend, and XLA settings. Changing them requires a new scope.

For the proposal-field preparation, four raw-coordinate starts $q_i$ and
positive parameter scales $s$ give the initial factor $A=\operatorname{diag}(s)$:
\begin{equation}
 c=\frac{1}{4}\sum_{i=1}^{4}q_i,
 \qquad z_i=(q_i-c)A^{-\mathsf T},
 \qquad q_i=c+z_iA^{\mathsf T}.
 \label{eq:bf-tensorflow-hmc-initial-bank}
\end{equation}
Here starts are represented as row vectors. Caller order is preserved.
A diagnostic may explicitly replicate one start; candidate preparation requires
its declared four-chain bank. Affine metric windows use $L=1$ and the configured
adaptation counts. Their proposals receive fresh fixed-kernel measurements.

The ordinary preparation config and numerical binding default to GPU/XLA
execution. TensorFlow memory growth is configured before device initialization.
Small CPU/non-XLA checks are explicit reference or debugging exceptions.
Execution policy does not qualify an arbitrary target: the binding still checks
its full-chain capability and the consumer must establish value/score fidelity.

\section{Broad search, repairs, and further evidence}
\label{sec:bf-fixed-transport-joint-tuning}
\label{sec:bf-fixed-transport-candidate-selection}

The ordinary convenience grid is $L\in\{3,5,9,13,18,25\}$, and its legacy
preparation config fixes the maximum at 25. The proposal-field config declares
its own maximum. Supplied initial, refinement and expansion grids must obey
the applicable maximum. An explicit \path{epsilon_by_l} may supply several proposals
at each $L$; \path{initial_epsilon} is a warm start when no grid is supplied.
Every pair is measured independently. A shared numeric proposal at two values
of $L$ transfers no acceptance evidence between them. The final-metric epsilon
from operational preparation supplies the starting step and the inherited search
ceiling. The short probe does not establish a stability limit or guarantee a
suitable pair below that ceiling for every $L$. An explicit
\path{candidate_search_bound_expansion_steps=1} widens the exploration cap by
the existing repair factor, preserving the final metric and coordinates.
Zero retains the inherited ceiling. Every explored pair still requires its own
health checks, measurement and fresh verification. This optional setting is a
search hypothesis whose calibration remains separate from its implementation.

Direct preparation bindings expose \path{preparation_bound_expansion_steps}.
They intersect the requested domain with the finite expanded cap, save the
original bound and final metric/coordinate identities, and issue a new search
identity. Existing scopes and verification receipts are immutable. Preparation
failures preserve their stage and available veto causes in
\path{preparation_progress.json} before raising an exception.
Nonfinite diagnostic values are marked explicitly in the JSON record, preserving
the original failure even when its numerical summaries cannot be finite.

A multiplicative repair can jump from high acceptance directly to nonfinite
trajectories. That nonfinite result supplies no valid acceptance direction,
so it cannot form the opposing-evidence interval used by refinement. The optional
\path{explore_failed_intervals=True} controller setting allows funded refinement
rounds to explore the unvisited geometric interior between a valid directional
parent and its nearest numerically rejected same-$L$ child. Each proposal is a
new same-family child of the valid parent; both source receipts are recorded,
and the failed endpoint stays rejected. This interval is an exploration
hypothesis, not a stability bound or an assertion of monotone acceptance.
Missing telemetry, unknown invalidity, invalid retained states and shared
corruption cannot supply such endpoints. Family, candidate and work limits
still apply, and every proposed pair needs measurement and independent
verification. The option is disabled by default. An explicit intermediate grid
in a fresh scope remains available. Neither procedure exhausts the continuous
epsilon domain.

The automatic preparation translation enables one pilot per $L$ and one bounded
refinement round. A direct controller config can disable these optional stages
for an explicitly supplied grid. A pilot proposes a frozen pair; measurement
and fresh verification qualify it. Refinement applies declared epsilon factors
and additional $L$ hypotheses to every surviving family. Exact pairs are
deduplicated, and all proposals remain subject to candidate and work limits.
The reciprocal convenience factors $(0.8,1.25)$ are hypotheses about useful
nearby proposals, not universal numerical optima.

Work proceeds in closed cohorts. Every funded admitted peer receives its reserved
stage before a repair child enters the next cohort. A peer whose full required
work cannot fit the remaining budget is explicitly deferred, permitting other
affordable reserved work to finish. Its evidence allocation is preserved and
the search remains incomplete. A first pass does not
remove later candidates. Directional evidence creates an immutable child with
the same $L$, mass, target, coordinates, and initial bank. The child has a new
epsilon and must receive its own measurement and verification. A parent may
fail movement because an oversized step rejects nearly everything; that
promotion veto does not suppress a supported smaller-epsilon repair.
A trajectory alert can propose only declared additional $L$ values, as separate
candidates. It never silently changes an existing candidate's $L$.

The distinction between proposal and evidence matters for fixed-length HMC.
For a transition from $x$, the acceptance probability is
\begin{equation}
 a(\epsilon,L\mid x)
 =\mathbb E\!\left[\min\{1,\exp(-\Delta H)\}\mid x\right].
\end{equation}
The reported probability mean averages the computed Metropolis probabilities
over visited states and sampled momenta. The realized acceptance rate also
contains the uniform accept/reject decisions: for each proposal with
probability $a$, that indicator has expectation $a$. Both summaries describe
acceptance in the explored region; neither measures convergence.
For the unit harmonic oscillator, one leapfrog matrix has determinant one and
trace $2-\epsilon^2$. In the stable range its conjugate eigenvalues therefore
give the phase
\begin{equation}
 \vartheta(\epsilon)=\arccos(1-\epsilon^2/2),\qquad 0<\epsilon<2.
\end{equation}
The trajectory phase is $L\vartheta(\epsilon)$. Near a return phase, acceptance
can be high with little net movement. A multiplicative epsilon repair is thus
another hypothesis to measure. The controller remembers the latest valid
directional evidence for each exact pair and prioritizes unvisited geometric
interior points toward nearby opposing observations. It retains this evidence
across restart. Finite refinement also investigates unresolved directional
intervals when no original family survives. It does not assume monotone
acceptance or infer a pass at an unmeasured point.
When a proposal crosses the declared domain, its unvisited boundary can be
measured before that direction is exhausted.

Inconclusive measurement or verification receives the predeclared
\path{evidence_rungs}; the convenience multipliers are $(1,2,4)$. Each rung is
an independently seeded run of the unchanged candidate. Unfinished work stays
resumable; exhaustion gives \path{inconclusive_at_cap}. The reported
compatibility intervals support an operational screen, not a claim of nominal
confidence coverage across repeated looks. A cap does not turn an inconclusive
candidate into a verified member.

\section{Acceptance, numerical health, and posterior diagnostics}

The default \path{HMCAcceptancePolicy} uses four independent chain means,
temporal blocks, and a 90\% compatibility interval. Its first allocation
requires at least four blocks of sixteen decisions per chain. After health,
chain-conflict, temporal-conflict and trajectory checks, acceptance passes when
the compatibility interval overlaps the practical band, lies inside the repair
band, and every chain mean lies inside the repair band. For example, mean
$0.76$ and interval $[0.73,0.79]$ can pass with practical band $[0.65,0.75]$
and repair band $[0.55,0.85]$. The rule tests compatibility; it does not rank
candidates by closeness to $0.70$. Acceptance,
movement and numerical health have distinct roles. Finite accepted/proposed
states and targets, endpoint score finiteness, momentum, Metropolis state
consistency and declared target status determine
whether the numerical evidence is usable. Candidate-local failure rejects that
pair; shared execution corruption stops the scope and disables all its replay.
A declared target-domain exception preserves the attempted seeds and failure
without inventing an acceptance value, so other candidates can continue.
Shared invalidity preserves the history of earlier verification while revoking
current replay. Released reservations cause affordable deferred work to be
reconsidered within the same invocation.
Available native divergence vetoes promotion of the current pair. When the
acceptance evidence is valid and supports a directional repair, that veto does
not prevent a newly measured child. Acceptance compatibility and promotion
eligibility are separate typed fields, so a rejected pair's in-band acceptance
receipt remains readable alongside its verified peers.
Both routes check numerical health over discarded warmup as well. The
proposal-field route checks finite states, log acceptance and energy errors
before excluding warmup from acceptance statistics.

$\widehat R$ is reporting-only during tuning. High, missing, nonfinite or
failed-to-compute values cannot reject, rank, repair or delay a candidate.
ESS and short-chain timing are descriptive diagnostics. They do not silently
replace the acceptance/health screen or discard a viable member.
Posterior warmup and cumulative model-coordinate $\widehat R$/ESS assessment
remain separate requirements for inference.

\section{A frozen nonlinear transport}

Let the base log density in model coordinates be
\(\ell_{\theta}(\theta)\), and let one frozen differentiable bijection map
latent coordinates into those model coordinates,
\[
  \theta = T(z), \qquad J_T(z)=\frac{\partial T(z)}{\partial z}.
\]
The fixed-transport HMC target is
\begin{equation}
  \ell_z(z)
  = \ell_{\theta}\!\left(T(z)\right)
    + \log\left|\det J_T(z)\right|,
  \label{eq:bf-fixed-transport-hmc-target}
\end{equation}
and its matching score is
\begin{equation}
  \nabla_z\ell_z(z)
  = J_T(z)^{\mathsf T}\nabla_{\theta}
      \ell_{\theta}\!\left(T(z)\right)
    + \nabla_z\log\left|\det J_T(z)\right|.
  \label{eq:bf-fixed-transport-hmc-score}
\end{equation}
Dropping either Jacobian term changes the target or its derivative.  Such a
calculation is wrong relative to the transformed-target claim.


The fixed-transport preparation reconstructs the full diagonal-affine or
dense-IAF map from \path{frozen_transport_payload}, checks its live identity,
and uses identity mass in latent coordinates. Starts are supplied in those
coordinates. Every proposed $(L,\epsilon)$ then enters the same measurement,
repair and fresh-verification procedure. Arbitrary transport callbacks without
a supported reconstruction are rejected. Changing a map changes the scope;
receipts or chain states from different maps are never pooled.

Positive funnel tuning tests assume a supplied frozen whitening or partially
whitening map. Learning a useful map is a separate preparation problem.
For the model $v\sim\mathcal N(0,\sigma^2)$ and
$x_i\mid v\sim\mathcal N(0,e^v)$ with two children, the exact map
$v=\sigma z_0$, $x_i=e^{v/2}z_i$ has log determinant
$\log\sigma+v$. Substitution in the transformed density cancels both
conditional scale terms and gives three independent standard normals.
This provides an algebraic control without assuming that a learner discovers
the map. A partial map $x_i=e^{s(v)}z_i$ instead leaves conditional standard
deviation $e^{v/2-s(v)}$ in the latent coordinates; it tests tuning with known
residual curvature.

The untransformed centered funnel remains a stress test for bounded search and
honest numerical rejection. An empty candidate set there does not by itself
establish a tuning defect. For a supplied map, tuning searches the declared
step sizes and trajectory lengths in the frozen coordinates and independently
verifies every survivor. A finite search need not succeed for every imperfect
map. Comparisons at common model starts must apply the inverse chart and
inverse map before passing latent starts to the tuner, then check the forward
roundtrip. Posterior precision and exploration in the original model remain
separate from candidate qualification.

\lstinputlisting[language=Python,caption={Fixed-transport candidate-set tuning.},label={lst:bf-hmc-tuning-fixed-transport}]{examples/hmc_tuning_fixed_transport.py}

\section{Typed deterministic proposal-field mechanics}

\path{bind_neural_force_hmc_tuning_runner} binds a declared position-only field
to its coordinate-consistent endpoint potential. The public facade accepts
that binding with \path{TensorFlowHMCKernelTuningConfig}. Affine preparation
remains specific to this transition, while the shared controller schedules its
candidate measurements, repairs and verification. The result is conditional
mechanics evidence and cannot become an exact-score numerical retained member.
The endpoint function still needs target-specific validation.

The typed objects use potential, not log-density, signs.  If
\(\ell(\theta)=\log \pi(\theta)\), the exact endpoint object is
\begin{equation}
  U(\theta)=-\ell(\theta).
  \label{eq:bf-neural-force-endpoint-potential}
\end{equation}
For the scalar-potential-gradient semantic, the proposal force is
\(F(\theta)=\nabla U(\theta)=-\nabla\ell(\theta)\).  The kick uses
\begin{equation}
  p_{n+1/2}=p_n-\frac{\epsilon}{2}F(q_n),
  \qquad
  q_{n+1}=q_n+\epsilon M^{-1}p_{n+1/2},
  \label{eq:bf-neural-force-kick}
\end{equation}
followed by the symmetric terminal kick and momentum flip.

The v2 binding also permits an honestly labeled deterministic position-only
proposal field \(B(q)\) that need not equal \(\nabla U(q)\).  Position-only
kicks remain volume-preserving shears, and the symmetric schedule remains
reversible.  The exact endpoint Metropolis rule is
\begin{equation}
  \alpha
  = \min\left\{1,
      \exp\bigl[-U(q')-K(p')+U(q)+K(p)\bigr]\right\}.
  \label{eq:bf-neural-force-endpoint-correction}
\end{equation}
Under these conditions, the endpoint potential determines the invariant target;
a poor or sign-reversed proposal field can destroy efficiency without by itself
changing that target.  A wrong endpoint potential, missing coordinate change,
momentum-dependent field, asymmetric schedule, or direct state update invalidates
that conclusion.

For an affine mass coordinate \(\theta=c+Fz\), the repository wrapper applies
the chain rule to the proposal field and endpoint potential.  The affine
log-Jacobian is constant and is deliberately omitted from energy differences
under the declared \code{constant\_omitted} convention.  Both raw objects and
that convention are bound into runner schema
\path{bayesfilter.hmc_tuning_runner_binding.v2}.


\lstinputlisting[language=Python,caption={Typed proposal field and exact endpoint-potential binding.},label={lst:bf-hmc-tuning-neural-force-binding}]{examples/hmc_tuning_neural_force_binding.py}

\section{Persistence, resource limits, and replay}
\label{sec:bf-hmc-replay-roles}

\path{verified_candidate_ids} retains every verified member with no implicit
nominee. \path{viable_candidate_ids} also includes validating and terminally
inconclusive candidates; their individual states must be inspected. A verified
member remains usable when another candidate exhausts its budget, unless the
scope itself is invalid. Descriptive differences among members do not establish
a statistically supported ranking.

A fresh output directory preserves \path{execution_spec.json}, immutable
numerical evidence and chunk files, and the atomic
\path{tuning_checkpoint.json}. All observations, including failed measurements
and zero-verified searches, are retained. Complete results use
\path{candidate_set_result.json} and are never overwritten. Resource failures
pause with their reason; unexpected exceptions preserve state and propagate.
Optional diagnostic errors cannot erase completed numerical evidence.
Malformed provider data is recorded as an execution error before it enters the
observation history; retry retains the attempted cost. Valid shared-failure
observations remain recorded. If the same live binding later records shared
invalidity, even a member verified earlier becomes unavailable for replay.

\path{resume_hmc_candidate_set_tuning} reconstructs the same geometry, policy,
evidence, candidates and work queue using the original target. It checks source
and target drift, tensor checksums, numerical versions, and scope identity.
Position-field mechanics separately preserve preparation and observations for
\path{resume_position_field_candidate_tuning}, which also requires the original
runner binding. Their \path{controller_checkpoint.json} also references immutable
completed chunks, permitting restart within a longer evidence stage.
Controller-only reload cannot create missing numerical evidence.

Before frozen execution exists, \path{preparation_progress.json} records
preparation phases, elapsed time and failures. Ordinary preparation checks time
at available progress boundaries; proposal-field affine preparation is one
compiled graph checked before and after its call. Failed preparation must be
retried in a fresh directory. Numerical resume requires a frozen scope with
the same source and controller policy; historical results remain readable when
those identities change, but cannot be silently resumed under new rules.

Call budgets count attempted work. Candidate reservations protect mandatory
stages; extensions and retries can use remaining free budget. A separate
transition-times-$(L+1)$ estimate accounts conservatively for gradient work.
The scheduler quotes unfinished work, and each chunk is charged and checkpointed
before its native call. Resume does not recharge saved completed chunks; an
interrupted native call remains conservatively charged even when it returned
no usable evidence. Elapsed time includes recorded preparation. Numerical chunks are bounded
by \path{chunk_max_results}, whose convenience default is 256. Time limits are
checked between chunks; Python cannot preempt a native operation or compilation.
A hard process limit must be applied externally. When preparation and search
time limits are both supplied, the smaller limit applies.
The first-call timing includes any tracing/compilation rather than pretending
that it measures steady-state kernel cost.

Select a member by explicit ID and supply its result and numerical binding to
\path{build_retained_bound_hmc_archive_runner_from_candidate_set_result}.
The two frozen-kernel variants share these checks, with the claim-eligible
variant additionally requiring actual GPU/XLA evidence. A child's own passing
verification is required; its failed ancestor need not pass. Export with
\path{runner.export(path)} and reconstruct with
\path{load_hmc_candidate_retained_runner(path, adapter=original_target)}.
Fresh retained blocks continue from the verified or preceding endpoint with
frozen settings and fresh seeds. Freshness checks include every tuning stage
and chunk, attempted partial chunks and predecessor retained blocks. Attempted
seeds are reconstructed from the durable work/accounting records after member
export and reload, including failed native calls that returned no trace.
Tuning draws remain excluded.
\path{runner.run_sequential} connects the explicitly chosen member to the
shared cumulative posterior controller, preserving the member's transition,
chain topology and target-status requirements. The generic entry point is
\path{run_hmc_posterior(member=runner, config=config)}. Diagnostics use model
coordinates and any declared scientific functionals. Posterior rejection does
not change tuning membership. Chapter~\ref{ch:bf-diagnostics} distinguishes
adaptation, discarded equilibration and retained precision.

The status \path{complete} means the declared search reached terminal
dispositions; its verified set may be empty. It does not prove the absence of a
useful kernel. The recorded limiting reason distinguishes the candidate cap,
the family repair count, the epsilon domain and exhaustion of representable
unvisited proposals. At a refinement cap, the remaining count includes only
eligible pairs not yet admitted in that round. Reporting collections resolve explicit scope/search/member
identities and require every included search to complete, independently of
their ordering.

Public option conflicts are rejected before numerical preparation. An issued
binding rejects redundant search, execution, provenance or transport overrides.
Automatic routes may accept the shared search and execution configs. For the
proposal-field route, legacy \path{step_adaptation_results} supplies the
fixed-pair pilot count, while \path{verification_results} supplies measurement
and verification counts. An explicit \path{HMCCandidateExecutionConfig} replaces
these allocations, with an optional \path{pilot_num_results}. Pilots never
qualify a pair. The legacy adaptive selector remains historical. Config
payloads label their legacy candidate-policy fields as historical helper
metadata; the active result's shared configs determine candidate stages.
Proposal-field execution defaults to XLA; a diagnostic exception supplies
\path{non_xla_reason}.

Retained-member exports share one checksummed evidence bundle beside their
compact member files. Copy the bundle when moving the members, or use
\path{runner.export(path, portable=True)} for a standalone file. Numerical
analyses are cached by current evidence content; mutations and changed source
identities still invalidate reuse. Predecessor history is checked iteratively. These mechanics tests do not qualify target-scale memory, restart
latency, posterior convergence, or scientific validity. Consumer evidence must
also establish complete target/data/prior identity and any intermediate
integrator telemetry required by that target.

\lstinputlisting[language=Python,caption={Ordinary candidate-set invocation.},label={lst:bf-hmc-tuning-ordinary}]{examples/hmc_tuning_ordinary.py}
\lstinputlisting[language=Python,caption={Numerical candidate evidence and explicit retained-member replay.}]{examples/hmc_candidate_set_retained.py}

\section{What the validation experiments establish}
\label{sec:bf-inference-validation}

The tuning procedure is tested through several distinct experiments. Numerical
oracles compare target values, scores, coordinate Jacobians and leapfrog updates
with independent calculations. Frozen-kernel experiments test invariance using
independent reference draws and a reversible random-position construction.
The latter uses Algorithm 2 and Proposition 2.3 of
\citet{gandy2021mcmctesting}; it applies to frozen reversible kernels.
An experiment may replace its kernel $K$ by $K^s$, composing $s$ complete
Metropolis transitions with independent randomness at each substep. Repeated
composition of this same reversible kernel preserves reversibility; section 2.2
of that paper proposes this extension to improve sensitivity to subtle errors.
Sensitivity still needs measurement for the chosen alternative. The option
\path{kernel_power} applies only to the frozen invariance experiment. It does
not make posterior draws independent. Finite states and log acceptance ratios are checked at every
substep, including earlier substeps when the final one is finite.

A separate optional wrapper implements Algorithm 3 of
\citet{gandy2021mcmctesting}. Each look uses an independently seeded complete
experiment, with the sample count increased once after the first look. For
overall level $\alpha$, at most $k$ looks and a fixed family of $d$ component
tests, initialize $\beta=\alpha/k$ and $\gamma=\beta^{1/k}$. At each look,
set $q=d\min_j p_j$. Reject when $q\leq\beta$, stop without rejection when
$q>\gamma+\beta$, and otherwise continue with $\beta\leftarrow\beta/\gamma$.
Theorem 3.1 bounds the probability of rejection under the null by $\alpha$
when the look vectors are independent and each component p-value is
superuniform. Fresh experiments are essential: repeatedly testing cumulative
posterior draws would not satisfy that independence condition. Numerical
failures invalidate the affected experiment; they do not count as defect
detection. This wrapper is a validation experiment and does not alter the
tuning acceptance screen, warmup checks or posterior precision stopping.

The Gaussian mechanics experiment also reconstructs the Metropolis log ratio
from an independent analytic density and the actual TFP endpoint momenta. It
checks that the returned state agrees with the recorded accept/reject decision.
This arithmetic check can expose a reversed energy correction even when a short
distributional experiment has little power to detect its effect.
For an identity-metric Gaussian, an additional fixed-look test uses the sum
of squared standardized endpoint coordinates. Starting from independent exact
Gaussian draws makes the endpoints independent under the correct frozen
kernel, so this sum has a chi-squared distribution with degrees of freedom
equal to the number of draws times the dimension. The experiment adjusts for
all declared radius, rank and energy tests. Adjacent states from one chain
cannot replace independent endpoints in this calculation. Null controls and
complete repeated experiments measure size and power for the declared defect;
one detected error does not establish broad sensitivity.
Controller tests check that every qualified pair is retained, each repair keeps
its own L, and fresh verification and restart preserve the numerical record.
Full-procedure experiments then call the public tuner and run the separate
posterior controller. Small integration cases reconstruct every verified member.
Larger replicated experiments may predeclare one member by tuning identity
before sampling; all other verified members remain recorded as unassessed for
that experiment.

Simulation-based calibration generates a parameter from a proper prior and a
dataset conditional on that parameter. If the fitted outputs are independent
exact posterior draws, the generating parameter's randomized rank among them
is uniform. This supplies the null distribution for the test. The implemented design uses independent
complete fits within each dataset and one fixed output per fit. It includes
data-dependent likelihood quantities because parameter ranks alone can miss a
procedure that ignores the observations. Failed fits remain in the denominator;
conditional ranks from successful fits cannot establish calibration for a
procedure that sometimes fails. Candidate siblings share a dataset and are not
independent calibration replications.
The rank argument follows \citet{talts2020sbc}; the likelihood quantities follow
the sensitivity analysis of \citet{modrak2023sbcquantities}.

The later M15 experiment obtained all 246 planned fitted outputs over 82 fresh
datasets, using three independently fitted outputs per dataset. Its rank tests
did not reject, but three rank draws give a coarse test. Independent analytic
controls in M16 showed that the 32-dataset normal design detected a shift of
one quarter of a posterior standard deviation in only 27 of 256 experiments,
and a half-standard-deviation shift in 64 of 256. The corresponding
nonrejection therefore supplies little evidence against those errors. These
controls measure the rank statistic's sensitivity; they do not measure defect
power of the complete HMC procedure. The full denominators, uncertainty and
source versions are recorded in the M15 and M16 result notes under
\path{docs/plans/}.

An invariant kernel can remain motionless or revisit only a few states.
Identity and two-cycle controls therefore accompany the distributional tests.
Similarly, an MCSE estimator can perform adequately on fixed-length stationary
arrays while intervals based on it have poor coverage at a data-dependent
stopping time. Separate experiments compare the actual controller's stopped
means and medians with available exact references, retaining cap and unavailable
outcomes. A fixed-length arm reuses the same verified kernel with independent
randomness and predeclared discarded and retained counts. Mixture mode
probabilities check global exploration that a local mean or small reported MCSE
may miss. Missing arms remain in the comparison denominator; pointwise
uncertainty intervals from a small number of replications do not establish a
ranking or nominal sequential coverage. No R-hat, ESS or MCSE assessment changes
tuning membership.
Coverage is reported both over all planned fits and conditional on an available
interval. Repeating a fit to one fixed dataset measures conditional behavior;
it is distinct from prior-predictive simulation-based calibration. A longer
fixed-count arm can also have inaccurate uncertainty estimates, so disagreement
with nominal coverage does not by itself identify optional stopping as the cause.

The September 2026 campaign observed this distinction directly. A Gaussian
mixture run passed its local posterior checks after 6,500 discarded and 2,000
retained transitions per chain, but all retained draws occupied the right mode.
Its estimated mean was 4.956, against the exact mean 2, with reported MCSE 0.038.
Thus a small within-mode uncertainty estimate did not measure error relative
to the full posterior. Target-specific posterior assessment should include
relevant mode or other global quantities and examine sensitivity to starting
locations. The saved result and its limited replication evidence are documented
in \path{docs/plans/bayesfilter-inference-validation-24h-result-2026-09-16.md}.
Validation designs may now declare the mixture probability $P(x<0)$ as a
monitored posterior mean. It receives the same diagnostic and precision checks
as the other declared quantities; a constant event indicator supplies no usable
precision estimate. A separate experiment starts chains in both known modes.
These checks address the saved failure mechanism without establishing discovery
of unknown modes. Synthetic acceptance calibration separately measures the
candidate screen near its practical and repair boundaries, including repeated
evidence allocations. It does not establish posterior validity.
The separate numerical acceptance experiment calls the public tuner on a
declared identity-metric Gaussian pair. Every repeated search measures that
exact pair from the fixed tuning starts, with no repair children. Each
measurement survivor receives its own fresh verification. Independent exact
Gaussian anchors supply a stationary acceptance
reference using actual Metropolis transitions and a separate endpoint-energy
check. That reference neither supplies tuning starts nor determines qualification.
Stationary acceptance and acceptance over a finite window from fixed starts
are different quantities; their difference is part of the experiment.
Observed qualification rates therefore describe the operational compatibility
screen, not nominal sequential coverage or automatic-preparation reliability.
The M16 experiment completed 384 CPU-reference and 192 GPU/XLA searches around
stationary acceptance values from .64 to .76. Pairs on both sides of the
nominal [.65,.75] band could qualify. This is consistent with the implemented
uncertainty-aware compatibility rule, whose finite-start evidence measures a
different quantity from the independent stationary reference. Separately, 128
complete sequential validation experiments per device detected the declared
reversed-Metropolis-ratio defect every time. Three of the four baseline/no-op
null intervals were still too wide for the declared size-precision screen.
That result supports detection of this defect while leaving null-size
precision and sensitivity to other defects unresolved.
Stopped-interval reporting compares the monitored mode probability with its
independent mixture CDF. That quantity remains in the planned denominator when
warmup ends without retained draws.

Automatic initialization and supplied epsilon proposals are distinct subjects
of validation. A supplied proposal can exceed the final prepared metric's bound;
native initialization uses preparation's own starting epsilon. Either search
can return no verified member under its finite domain and budget. A complete
calibration experiment must preserve those missing fits, even when all observed
ranks among successful fits look uniform. Parallel dataset groups are combined
only when declared together with identical settings and frozen source; pilot and
repaired experiments remain separate.

Matched consumer comparisons use the same public ordinary tuner and posterior
controller. The explicit noncentered eight-schools and Bayesian regression
adapters reproduce the pinned posteriordb Stan laws, observations and coordinate
Jacobians. A member is selected by tuning identity before posterior sampling;
every other verified member remains recorded. The ten-chain Stan reference is
used only for the subsequent comparison, with lugsail uncertainty from both
finite samples. Agreement within a declared practical margin is a
target-specific accuracy screen. It neither makes the reference exact nor
establishes calibration for another dataset or consumer. Exact inputs, commands
and results are recorded in
\path{docs/plans/bayesfilter-hmc-repair-m8-result-2026-09-18.md}.

The later M17 matrix completed 21 model/route fits and four matched-reference
fits. Eleven matrix members and three reference members passed their full
declared posterior checks. Centered funnels produced no verified member;
mixtures reached the warmup cap; rotated Gaussian, noncentered funnel and CPU
eight-schools fits reached retained precision caps. Their verified kernels
remain in the tuning set. All four reference fits met the declared mean
agreement margins, illustrating why reference mean agreement and requested
precision must be reported separately. Six conditional position-field cases
also exercised fresh verification and durable restart. Fixed nonlinear maps
in this matrix were not trained. Exact model/device outcomes and source versions
are recorded in \path{docs/plans/bayesfilter-hmc-repair-m17-result-2026-09-21.md}.
The original MacroFinance target and matched posterior reference remain a
separate missing comparison.

Cost pilots can save a host profile of automatic preparation, candidate search,
export/reload and posterior assessment. The profile includes failed attempts;
a framework call may include both compilation and execution time. Such a pilot
identifies costs to investigate before funding more replications. It does not
establish statistical calibration or rank the viable kernels by efficiency.

Repeated validation fits can use a fresh process for each complete fit, with
an explicit per-fit timeout inside the total experiment budget. This bounds
the lifetime of process-local framework caches. Each child preserves its full
candidate set, posterior chunks and independent assessment before normal exit;
the coordinator records process failures separately from saved numerical
results. Missing or abnormally terminated fits remain in the replication
denominator. Numerical search, accuracy and stopping designs enable this
optional mode with \path{options.isolate_fits} and
\path{options.fit_process_timeout_seconds}. The setting does not change the
HMC transition, tuning criteria or posterior thresholds.

A bounded sibling study can instead predeclare several members using
\path{member_rule="shortest_verified_l"},
\path{posterior_members="selected"} and an explicit
\path{posterior_member_count}. The rule orders distinct verified leapfrog
counts increasingly and takes the first candidate identifier within each.
It records the list before posterior sampling and preserves any shortage.
Each ordinal member is assessed separately across complete fits; siblings do
not increase the number of independent replications. A failed member does not
suppress the next member or remove a tuning candidate. This is a bounded cost
hypothesis, not an established mixing ranking, and posterior outcomes cannot
choose a winner from the same assessment draws.

The executable suite definitions are documented in
\path{docs/validation/README.md}. The following table is generated from their
saved reports. C, H and P denote current-source executions, historical-source
executions and planned designs, including retries. These are counts of designs,
not models declared correct or independent replications. CPU reference work and
GPU/XLA work remain separate. A completed controller call with no retained draws
does not establish posterior coverage; missing stopping intervals remain in the
denominator. Affine and dense nonlinear transports occupy separate rows.
Each experiment declares its evidence scope. Incomplete confirmation cannot
establish calibration; statistical uncertainty and unresolved coverage are
reported with the underlying observations.

The table covers suite reports. Separately executed phase diagnostics, including
the matched posteriordb comparisons above, are recorded in their phase results.
An unavailable external eight-schools row refers to that original suite's
missing reference bundle; it does not negate the later explicitly bound
eight-schools comparison. Neither supplies the absent MacroFinance reference.

Automatically inferred suite requirements identify each planned design, so a
completed step-size or kernel-power setting cannot cover an unrun setting in
the same category. Explicit category requirements may request any matching
assessed design. Older saved plans without design identifiers retain that
category-level meaning; their individual rows determine whole-suite completion.

% Generated by scripts/render_inference_validation_coverage.py
\begingroup\small\setlength{\tabcolsep}{3pt}
\begin{longtable}{>{\raggedright\arraybackslash}p{.15\textwidth}>{\raggedright\arraybackslash}p{.19\textwidth}>{\raggedright\arraybackslash}p{.13\textwidth}>{\raggedright\arraybackslash}p{.14\textwidth}>{\raggedright\arraybackslash}p{.28\textwidth}}
Target & Route/device & Question & C / H / P & Findings \\\hline\endhead
Banana & Frozen kernel / CPU & Invariance & 0 / 2 / 2 & no discrepancy detected \\
Banana & Frozen kernel / CPU & Mechanics & 0 / 3 / 3 & mechanics passed \\
Banana & Frozen kernel / GPU & Mechanics & 0 / 1 / 2 & mechanics passed, not run \\
Beta & Frozen kernel / CPU & Invariance & 0 / 2 / 2 & no discrepancy detected \\
Beta & Frozen kernel / CPU & Mechanics & 0 / 3 / 3 & mechanics passed \\
Beta & Frozen kernel / GPU & Mechanics & 0 / 1 / 2 & mechanics passed, not run \\
Beta & Prepared / CPU & Accuracy & 0 / 1 / 1 & pipeline assessed \\
Beta-binomial & Frozen kernel / CPU & Mechanics & 0 / 3 / 3 & mechanics passed \\
Beta-binomial & Frozen kernel / GPU & Mechanics & 0 / 1 / 2 & mechanics passed, not run \\
Beta-binomial & Automatic / CPU & SBC & 0 / 1 / 1 & calibration incomplete \\
Beta-binomial & Automatic / GPU & SBC & 0 / 1 / 2 & no discrepancy detected, not run \\
Beta-binomial & Reference / CPU & SBC & 0 / 2 / 2 & no discrepancy detected \\
Cauchy & Frozen kernel / CPU & Invariance & 0 / 2 / 2 & no discrepancy detected \\
Cauchy & Frozen kernel / CPU & Mechanics & 0 / 3 / 3 & mechanics passed \\
Cauchy & Frozen kernel / GPU & Mechanics & 0 / 1 / 2 & mechanics passed, not run \\
Cauchy & Automatic / CPU & Accuracy & 0 / 1 / 1 & pipeline assessed \\
Dirichlet & Frozen kernel / CPU & Invariance & 0 / 2 / 2 & no discrepancy detected \\
Dirichlet & Frozen kernel / CPU & Mechanics & 0 / 3 / 3 & mechanics passed \\
Dirichlet & Frozen kernel / GPU & Mechanics & 0 / 1 / 2 & mechanics passed, not run \\
Dirichlet & Prepared / CPU & Accuracy & 0 / 1 / 1 & pipeline assessed \\
Dirichlet & Prepared / GPU & Accuracy & 0 / 1 / 3 & not run, pipeline assessed, timed out \\
Eight schools & External / CPU & Accuracy & 0 / 0 / 2 & unavailable \\
Funnel & Frozen kernel / CPU & Invariance & 0 / 2 / 2 & no discrepancy detected \\
Funnel & Frozen kernel / CPU & Mechanics & 0 / 3 / 3 & mechanics passed \\
Funnel & Frozen kernel / GPU & Mechanics & 0 / 2 / 3 & mechanics passed, not run \\
Funnel & Automatic / GPU & Accuracy & 0 / 0 / 2 & failed, not run \\
Gamma & Frozen kernel / CPU & Invariance & 0 / 2 / 2 & no discrepancy detected \\
Gamma & Frozen kernel / CPU & Mechanics & 0 / 3 / 3 & mechanics passed \\
Gamma & Frozen kernel / GPU & Mechanics & 0 / 2 / 3 & mechanics passed, not run \\
Gamma & Prepared / CPU & Accuracy & 0 / 1 / 1 & pipeline assessed \\
Gaussian & Controller / CPU & Search & 0 / 2 / 2 & inventory passed \\
Gaussian & Dense IAF / GPU & Accuracy & 0 / 3 / 3 & no verified members, pipeline assessed \\
Gaussian & Affine transport / CPU & Accuracy & 0 / 1 / 1 & pipeline assessed \\
Gaussian & Affine transport / GPU & Accuracy & 0 / 1 / 1 & pipeline assessed \\
Gaussian & Frozen kernel / CPU & Invariance & 0 / 2 / 2 & no discrepancy detected \\
Gaussian & Frozen kernel / CPU & Mechanics & 0 / 3 / 3 & mechanics passed \\
Gaussian & Frozen kernel / CPU & Power & 0 / 3 / 3 & power estimated \\
Gaussian & Frozen kernel / GPU & Invariance & 0 / 1 / 1 & no discrepancy detected \\
Gaussian & Frozen kernel / GPU & Mechanics & 0 / 3 / 4 & mechanics passed, not run \\
Gaussian & Frozen kernel / GPU & Power & 0 / 4 / 7 & not run, power estimated \\
Gaussian & Automatic / CPU & Accuracy & 0 / 1 / 2 & pipeline assessed, timed out \\
Gaussian & Automatic / GPU & Accuracy & 0 / 1 / 2 & pipeline assessed, timed out \\
Gaussian & Prepared / CPU & Accuracy & 0 / 1 / 2 & pipeline assessed, timed out \\
Gaussian & Prepared / CPU & Stopping & 0 / 4 / 5 & pipeline assessed, timed out \\
Gaussian & Prepared / GPU & Stopping & 0 / 3 / 5 & not run, pipeline assessed \\
Gaussian & Reference / CPU & Stopping & 0 / 6 / 6 & diagnostics assessed \\
LGSSM location & Frozen kernel / CPU & Mechanics & 0 / 3 / 3 & mechanics passed \\
LGSSM location & Frozen kernel / GPU & Mechanics & 0 / 1 / 2 & mechanics passed, not run \\
LGSSM location & Automatic / CPU & SBC & 0 / 1 / 1 & no discrepancy detected \\
LGSSM location & Automatic / GPU & SBC & 0 / 1 / 2 & calibration incomplete, not run \\
LGSSM location & Reference / CPU & SBC & 0 / 2 / 2 & no discrepancy detected \\
Macrofinance & External / CPU & Accuracy & 0 / 0 / 2 & unavailable \\
Mixture & Frozen kernel / CPU & Invariance & 0 / 2 / 2 & no discrepancy detected \\
Mixture & Frozen kernel / CPU & Mechanics & 0 / 3 / 3 & mechanics passed \\
Mixture & Frozen kernel / GPU & Mechanics & 0 / 1 / 2 & mechanics passed, not run \\
Mixture & Prepared / CPU & Stopping & 0 / 3 / 4 & incomplete, pipeline discrepancy, timed out \\
Mixture & Prepared / GPU & Stopping & 0 / 2 / 4 & not run, pipeline discrepancy \\
Mixture & Reference / CPU & Stopping & 0 / 2 / 2 & diagnostics assessed \\
Normal-normal & Frozen kernel / CPU & Mechanics & 0 / 3 / 3 & mechanics passed \\
Normal-normal & Frozen kernel / GPU & Mechanics & 0 / 1 / 2 & mechanics passed, not run \\
Normal-normal & Automatic / CPU & SBC & 0 / 22 / 22 & calibration incomplete, no discrepancy detected \\
Normal-normal & Automatic / GPU & SBC & 0 / 4 / 15 & calibration incomplete, cancelled, not run, timed out \\
Normal-normal & Prepared / CPU & SBC & 0 / 17 / 17 & no discrepancy detected \\
Normal-normal & Reference / CPU & Power & 0 / 9 / 9 & power estimated \\
Normal-normal & Reference / CPU & SBC & 0 / 2 / 2 & no discrepancy detected \\
Regression & External / CPU & Accuracy & 0 / 0 / 2 & unavailable \\
Rotated gaussian & Affine transport / CPU & Accuracy & 0 / 1 / 1 & posterior incomplete \\
Rotated gaussian & Affine transport / GPU & Accuracy & 0 / 1 / 2 & not run, posterior incomplete \\
Rotated gaussian & Frozen kernel / CPU & Invariance & 0 / 2 / 2 & no discrepancy detected \\
Rotated gaussian & Frozen kernel / CPU & Mechanics & 0 / 3 / 3 & mechanics passed \\
Rotated gaussian & Frozen kernel / GPU & Mechanics & 0 / 1 / 2 & mechanics passed, not run \\
Student t & Frozen kernel / CPU & Invariance & 0 / 2 / 2 & no discrepancy detected \\
Student t & Frozen kernel / CPU & Mechanics & 0 / 3 / 3 & mechanics passed \\
Student t & Frozen kernel / GPU & Mechanics & 0 / 1 / 2 & mechanics passed, not run \\
\end{longtable}
\endgroup


\section{Historical readers and complete interface inventory}

Initial geometry and bootstrap screening are implemented in
\path{hmc_geometry.py} and \path{hmc_bootstrap.py}, respectively.
Windowed preparation, its timeout policy and the frozen-mass/start-bank handoff
are implemented in \path{hmc_mass_adaptation.py}.
Public imports and historical \path{hmc_kernel_tuning} aliases resolve to the
same definitions. Automatic preparation calls these implementations directly;
public presets, configuration translation and geometry-scaled budgets live in
\path{hmc_configuration.py}. The old module retains historical orchestration and
readers, with aliases for the extracted definitions.
These helpers do not issue candidate-set tuning authority. Numerical bindings
include their source files and shared preparation helpers. Seed identifiers
remain stable, while live seed records name the current implementation files.
Historical results retain their original source provenance.
\label{sec:bf-hmc-complete-interface-inventory}

Earlier single-kernel ordinary, fixed-transport and typed TensorFlow results
retain their original identities and readers. Their compatibility replay roles
and older \path{claim_bearing_blockers} remain historical evidence; a
caller-edited authority flag or blocker list cannot grant authority. The active
procedure uses the candidate-set result and the member's own numerical evidence.
Discovery, refinement, \path{select_fixed_transport_candidate_set}, and stage
helpers remain diagnostics. The generated registry below identifies their
replacement public route and their limits.

% Generated by scripts/render_hmc_tuning_interface_docs.py. Do not edit.
\begingroup
\footnotesize
\sloppy
\setlength{\tabcolsep}{3pt}
\renewcommand{\arraystretch}{1.15}
\begin{longtable}{>{\raggedright\arraybackslash}p{0.26\textwidth}>{\raggedright\arraybackslash}p{0.14\textwidth}>{\raggedright\arraybackslash}p{0.53\textwidth}}
\toprule
Interface & Status & Target, tuning, verification, and use \\
\midrule
\endfirsthead
\toprule
Interface & Status & Target, tuning, verification, and use \\
\midrule
\endhead
\path{bayesfilter.inference.hmc_tuning_dispatch.tune_hmc_kernel} & public tuner; tested supported; authority & Target: artifact authority requires an ordinary exact value/score target; the typed position-field branch additionally requires a repository-issued binding and exact endpoint target but remains mechanics-only. Coordinates: ordinary adapter coordinates with repository-owned affine mass coordinates. Mass: windowed adaptation by default or explicit fixed identity from config. Epsilon: shared per-L pilot, exact-pair measurement, bounded directional repair, and independently seeded evidence rungs. L: shared primary L grid (3, 5, 9, 13, 18, 25), explicit optional all-survivor refinement and expansion; retain every verified pair. Verification: fresh fixed-kernel verification; default TFP runner requires typed acceptance, health, and minimum draws; rank-normalized split/folded R-hat is retained as an explanatory diagnostic with threshold metadata (1.01), not a tuning handoff gate. ESS: disabled for ordinary tuning admission; retained posterior ESS is separate. Forbidden: arbitrary bare runner callback; engineering\_probe\_covariance\_multiplier as a public ordinary mode; transferring epsilon qualification across different L candidates; fixed nonlinear transport without its transformed target contract; retained posterior convergence claim. \\
\addlinespace
\path{bayesfilter.inference.fixed_transport_hmc_tuning_tf.tune_fixed_transport_hmc_kernel} & public tuner; tested supported; authority & Target: identity-bound frozen transport with exact transformed value including Jacobian and matching score. Coordinates: fixed transport latent z coordinates. Mass: fixed identity mass in z; no ordinary windowed mass adaptation. Epsilon: shared per-L pilot or explicit epsilon proposals, exact-pair measurement and bounded independently verified repairs. L: all declared (epsilon, L) pairs measured; all-survivor refinement and retention. Verification: fresh fixed-kernel verification for every survivor; declared evidence rungs; R-hat reporting-only. ESS: disabled for kernel tuning admission; cumulative posterior ESS is separate. Forbidden: arbitrary position-only force without a frozen transport; ordinary mass adaptation claim; reuse after transport identity changes; legacy directional diagnostic policy as an artifact-authority handoff. \\
\addlinespace
\path{bayesfilter.inference.hmc_candidate_set_tuning.tune_hmc_candidate_set} & diagnostic helper; diagnostic only; no authority & Target: not authoritative; use the active replacement target contract. Coordinates: not authoritative; use the active replacement coordinates. Mass: not authoritative for canonical handoff. Epsilon: not authoritative for canonical handoff. L: not authoritative for canonical handoff. Verification: cannot issue canonical fresh-verification admission. ESS: not authoritative for canonical handoff. Forbidden: canonical tuning artifact; retained-kernel handoff; public tuning recommendation. Replacement: tune\_hmc\_kernel. \\
\addlinespace
\path{bayesfilter.inference.fixed_transport_candidate_selection.select_fixed_transport_candidate_set} & diagnostic helper; diagnostic only; no authority & Target: not authoritative; use the active replacement target contract. Coordinates: not authoritative; use the active replacement coordinates. Mass: not authoritative for canonical handoff. Epsilon: not authoritative for canonical handoff. L: not authoritative for canonical handoff. Verification: cannot issue canonical fresh-verification admission. ESS: not authoritative for canonical handoff. Forbidden: canonical tuning artifact; retained-kernel handoff; public tuning recommendation. Replacement: tune\_hmc\_kernel. \\
\addlinespace
\path{bayesfilter.inference.hmc_robust_broad_grid.tune_hmc_kernel_robust_broad_grid} & diagnostic helper; diagnostic only; no authority & Target: not authoritative; use the active replacement target contract. Coordinates: not authoritative; use the active replacement coordinates. Mass: not authoritative for canonical handoff. Epsilon: not authoritative for canonical handoff. L: not authoritative for canonical handoff. Verification: cannot issue canonical fresh-verification admission. ESS: not authoritative for canonical handoff. Forbidden: canonical tuning artifact; retained-kernel handoff; public tuning recommendation. Replacement: tune\_hmc\_kernel. \\
\addlinespace
\path{bayesfilter.inference.hmc_fixed_metric_grid_search.run_fixed_metric_grid_search} & diagnostic helper; diagnostic only; no authority & Target: not authoritative; use the active replacement target contract. Coordinates: not authoritative; use the active replacement coordinates. Mass: not authoritative for canonical handoff. Epsilon: not authoritative for canonical handoff. L: not authoritative for canonical handoff. Verification: cannot issue canonical fresh-verification admission. ESS: not authoritative for canonical handoff. Forbidden: canonical tuning artifact; retained-kernel handoff; public tuning recommendation. Replacement: tune\_hmc\_kernel. \\
\addlinespace
\path{bayesfilter.inference.hmc_operational_broad_grid.run_operational_broad_grid} & diagnostic helper; diagnostic only; no authority & Target: not authoritative; use the active replacement target contract. Coordinates: not authoritative; use the active replacement coordinates. Mass: not authoritative for canonical handoff. Epsilon: not authoritative for canonical handoff. L: not authoritative for canonical handoff. Verification: cannot issue canonical fresh-verification admission. ESS: not authoritative for canonical handoff. Forbidden: canonical tuning artifact; retained-kernel handoff; public tuning recommendation. Replacement: tune\_hmc\_kernel. \\
\addlinespace
\path{bayesfilter.inference.hmc_operational_broad_grid.run_operational_broad_grid_process_parallel} & diagnostic helper; diagnostic only; no authority & Target: not authoritative; use the active replacement target contract. Coordinates: not authoritative; use the active replacement coordinates. Mass: not authoritative for canonical handoff. Epsilon: not authoritative for canonical handoff. L: not authoritative for canonical handoff. Verification: cannot issue canonical fresh-verification admission. ESS: not authoritative for canonical handoff. Forbidden: canonical tuning artifact; retained-kernel handoff; public tuning recommendation. Replacement: tune\_hmc\_kernel. \\
\addlinespace
\path{bayesfilter.inference.hmc_budget_ladder.run_fixed_mass_hmc_tuning_budget_ladder} & diagnostic helper; diagnostic only; no authority & Target: not authoritative; use the active replacement target contract. Coordinates: not authoritative; use the active replacement coordinates. Mass: not authoritative for canonical handoff. Epsilon: not authoritative for canonical handoff. L: not authoritative for canonical handoff. Verification: cannot issue canonical fresh-verification admission. ESS: not authoritative for canonical handoff. Forbidden: canonical tuning artifact; retained-kernel handoff; public tuning recommendation. Replacement: tune\_hmc\_kernel. \\
\addlinespace
\path{bayesfilter.inference.hmc_tuning.run_fixed_mass_step_tuning_diagnostic} & diagnostic helper; diagnostic only; no authority & Target: not authoritative; use the active replacement target contract. Coordinates: not authoritative; use the active replacement coordinates. Mass: not authoritative for canonical handoff. Epsilon: not authoritative for canonical handoff. L: not authoritative for canonical handoff. Verification: cannot issue canonical fresh-verification admission. ESS: not authoritative for canonical handoff. Forbidden: canonical tuning artifact; retained-kernel handoff; public tuning recommendation. Replacement: tune\_hmc\_kernel. \\
\addlinespace
\path{bayesfilter.inference.hmc_tuning.run_windowed_mass_adaptation_diagnostic} & diagnostic helper; diagnostic only; no authority & Target: not authoritative; use the active replacement target contract. Coordinates: not authoritative; use the active replacement coordinates. Mass: not authoritative for canonical handoff. Epsilon: not authoritative for canonical handoff. L: not authoritative for canonical handoff. Verification: cannot issue canonical fresh-verification admission. ESS: not authoritative for canonical handoff. Forbidden: canonical tuning artifact; retained-kernel handoff; public tuning recommendation. Replacement: tune\_hmc\_kernel. \\
\addlinespace
\path{bayesfilter.inference.hmc_tuning.run_fixed_trajectory_tuning_diagnostic} & diagnostic helper; diagnostic only; no authority & Target: not authoritative; use the active replacement target contract. Coordinates: not authoritative; use the active replacement coordinates. Mass: not authoritative for canonical handoff. Epsilon: not authoritative for canonical handoff. L: not authoritative for canonical handoff. Verification: cannot issue canonical fresh-verification admission. ESS: not authoritative for canonical handoff. Forbidden: canonical tuning artifact; retained-kernel handoff; public tuning recommendation. Replacement: tune\_hmc\_kernel. \\
\addlinespace
\path{bayesfilter.inference.hmc_tuning.run_gaussian_dual_averaging_diagnostic} & diagnostic helper; diagnostic only; no authority & Target: not authoritative; use the active replacement target contract. Coordinates: not authoritative; use the active replacement coordinates. Mass: not authoritative for canonical handoff. Epsilon: not authoritative for canonical handoff. L: not authoritative for canonical handoff. Verification: cannot issue canonical fresh-verification admission. ESS: not authoritative for canonical handoff. Forbidden: canonical tuning artifact; retained-kernel handoff; public tuning recommendation. Replacement: tune\_hmc\_kernel. \\
\addlinespace
\path{bayesfilter.inference.hmc_kernel_tuning.run_hmc_start_bank_diagnostic} & diagnostic helper; diagnostic only; no authority & Target: not authoritative; use the active replacement target contract. Coordinates: not authoritative; use the active replacement coordinates. Mass: not authoritative for canonical handoff. Epsilon: not authoritative for canonical handoff. L: not authoritative for canonical handoff. Verification: cannot issue canonical fresh-verification admission. ESS: not authoritative for canonical handoff. Forbidden: canonical tuning artifact; retained-kernel handoff; public tuning recommendation. Replacement: tune\_hmc\_kernel. \\
\addlinespace
\path{bayesfilter.inference.fixed_transport_hmc_candidate_discovery_tf.discover_fixed_transport_hmc_candidates} & diagnostic helper; diagnostic only; no authority & Target: not authoritative; use the active replacement target contract. Coordinates: not authoritative; use the active replacement coordinates. Mass: not authoritative for canonical handoff. Epsilon: not authoritative for canonical handoff. L: not authoritative for canonical handoff. Verification: cannot issue canonical fresh-verification admission. ESS: not authoritative for canonical handoff. Forbidden: canonical tuning artifact; retained-kernel handoff; public tuning recommendation. Replacement: tune\_fixed\_transport\_hmc\_kernel. \\
\addlinespace
\path{bayesfilter.inference.fixed_transport_hmc_candidate_discovery_tf.run_fixed_transport_hmc_candidate_campaign} & diagnostic helper; diagnostic only; no authority & Target: not authoritative; use the active replacement target contract. Coordinates: not authoritative; use the active replacement coordinates. Mass: not authoritative for canonical handoff. Epsilon: not authoritative for canonical handoff. L: not authoritative for canonical handoff. Verification: cannot issue canonical fresh-verification admission. ESS: not authoritative for canonical handoff. Forbidden: canonical tuning artifact; retained-kernel handoff; public tuning recommendation. Replacement: tune\_fixed\_transport\_hmc\_kernel. \\
\addlinespace
\path{bayesfilter.inference.fixed_transport_hmc_candidate_discovery_tf.refine_fixed_transport_hmc_candidates} & diagnostic helper; diagnostic only; no authority & Target: not authoritative; use the active replacement target contract. Coordinates: not authoritative; use the active replacement coordinates. Mass: not authoritative for canonical handoff. Epsilon: not authoritative for canonical handoff. L: not authoritative for canonical handoff. Verification: cannot issue canonical fresh-verification admission. ESS: not authoritative for canonical handoff. Forbidden: canonical tuning artifact; retained-kernel handoff; public tuning recommendation. Replacement: tune\_fixed\_transport\_hmc\_kernel. \\
\addlinespace
\path{bayesfilter.inference.generic_hmc_tuning.run_generic_hmc_tuning_orchestration} & historical helper; historical only; no authority & Target: not authoritative; use the active replacement target contract. Coordinates: not authoritative; use the active replacement coordinates. Mass: not authoritative for canonical handoff. Epsilon: not authoritative for canonical handoff. L: not authoritative for canonical handoff. Verification: cannot issue canonical fresh-verification admission. ESS: not authoritative for canonical handoff. Forbidden: canonical tuning artifact; retained-kernel handoff; public tuning recommendation. Replacement: tune\_hmc\_kernel. \\
\addlinespace
\path{bayesfilter.inference.generic_hmc_tuning.orchestrate_generic_hmc_tuning} & historical helper; historical only; no authority & Target: not authoritative; use the active replacement target contract. Coordinates: not authoritative; use the active replacement coordinates. Mass: not authoritative for canonical handoff. Epsilon: not authoritative for canonical handoff. L: not authoritative for canonical handoff. Verification: cannot issue canonical fresh-verification admission. ESS: not authoritative for canonical handoff. Forbidden: canonical tuning artifact; retained-kernel handoff; public tuning recommendation. Replacement: tune\_hmc\_kernel. \\
\addlinespace
\path{bayesfilter.inference.fixed_trajectory_hmc_tuning_v2.run_tiny_gaussian_fixed_trajectory_hmc_tuning_v2} & historical helper; historical only; no authority & Target: not authoritative; use the active replacement target contract. Coordinates: not authoritative; use the active replacement coordinates. Mass: not authoritative for canonical handoff. Epsilon: not authoritative for canonical handoff. L: not authoritative for canonical handoff. Verification: cannot issue canonical fresh-verification admission. ESS: not authoritative for canonical handoff. Forbidden: canonical tuning artifact; retained-kernel handoff; public tuning recommendation. Replacement: tune\_hmc\_kernel. \\
\addlinespace
\path{bayesfilter.inference.hmc_kernel_tuning.run_hmc_tune_verify_repair_loop} & stage helper; internal only; no authority & Target: prevalidated ordinary adapter, geometry, and bootstrap artifacts. Coordinates: same adapter and mass-coordinate lineage as public tuner. Mass: internal windowed adaptation and repair. Epsilon: internal fixed-mass tuning and repair. L: internal leapfrog-count selection and repair. Verification: internal fresh fixed-kernel verification. ESS: disabled for ordinary tuning admission. Forbidden: final one-call API; canonical artifact authority. Replacement: tune\_hmc\_kernel. \\
\addlinespace
\path{bayesfilter.inference.hmc.run_full_chain_tfp_hmc} & chain runner; internal only; no authority & Target: adapter exact value and matching score. Coordinates: caller-supplied fixed mass coordinates. Mass: not owned; consumes supplied coordinates. Epsilon: fixed or config-requested dual averaging only. L: fixed caller-supplied leapfrog count. Verification: none by itself. ESS: none by itself. Forbidden: standalone full-tuning claim; canonical artifact authority. Replacement: tune\_hmc\_kernel. \\
\addlinespace
\path{bayesfilter.inference.fixed_transport_hmc_mechanics_tf.run_fixed_transport_full_chain_tfp_hmc} & chain runner; internal only; no authority & Target: preconstructed exact transformed value/score adapter. Coordinates: frozen-transport latent z coordinates. Mass: not owned; consumes the supplied fixed coordinate system. Epsilon: fixed or config-requested dual averaging only. L: fixed caller-supplied leapfrog count. Verification: none by itself. ESS: none by itself. Forbidden: standalone full-tuning claim; candidate-selection claim; canonical artifact authority. Replacement: tune\_fixed\_transport\_hmc\_kernel. \\
\addlinespace
\path{bayesfilter.inference.neural_force_hmc.run_full_chain_neural_force_hmc} & chain runner; diagnostic only; no authority & Target: frozen position-only force plus exact deterministic endpoint potential. Coordinates: native affine mass coordinates when bound; direct identity fallback is diagnostic only. Mass: not owned; identity in supplied mass coordinates. Epsilon: fixed or config-requested dual averaging only. L: fixed caller-supplied leapfrog count. Verification: none by itself. ESS: none by itself. Forbidden: standalone full-tuning claim; mass-tuning claim; leapfrog-count-tuning claim; canonical artifact authority. Replacement: tune\_hmc\_kernel. \\
\addlinespace
\path{bayesfilter.inference.neural_force_hmc.bind_neural_force_hmc_tuning_runner} & runner binding factory; tested supported; no authority & Target: coordinate-consistent frozen force and exact raw-coordinate endpoint potential. Coordinates: TensorFlow tuning branch native affine mass wrapper; identity fallback for a direct low-level call is diagnostic only. Mass: owned by the typed TensorFlow mechanics branch. Epsilon: owned by the typed TensorFlow mechanics branch. L: powers-of-two candidate screen owned by the typed TensorFlow mechanics branch; not the ordinary broad-grid policy. Verification: fresh injected-runner fixed-kernel health and acceptance verification. ESS: disabled for ordinary tuning admission. Forbidden: standalone tuning authority; bare callable injection; transformed target masquerading as ordinary raw coordinates. Replacement: tune\_hmc\_kernel. \\
\addlinespace
\bottomrule
\end{longtable}
\endgroup

